%
% Documento: Introdução
\setcounter{page}{6}
\chapter{INTRODUÇÃO}\label{chap:introducao}


A educação inclusiva é um dos desafios enfrentados pelo sistema educacional brasileiro, principalmente quando se trata de garantir condições adequadas de aprendizagem para estudantes com deficiência intelectual e outras necessidades educacionais específicas. Embora a legislação brasileira assegure o direito à educação inclusiva, colocando esse direito em prática exige mais do que a matrícula dos estudantes nas escolas. A Lei Brasileira de Inclusão da Pessoa com Deficiência, instituída pela Lei nº 13.146/2015, estabelece diretrizes para garantir o acesso, a permanência, a participação e a aprendizagem das pessoas com deficiência no ambiente escolar \cite{brasil2015lbi}. No entanto, é necessário oferecer materiais pedagógicos adaptados, estratégias de ensino adequadas e acompanhamento individualizado, respeitando as características e necessidades de cada aluno.

Mesmo com os avanços na legislação, muitas escolas públicas ainda enfrentam dificuldades para oferecer um ensino inclusivo. A falta de materiais adaptados, a formação insuficiente dos professores, a escassez de profissionais especializados e a sobrecarga de trabalho são alguns dos obstáculos encontrados no ambiente escolar. \citeonline{souza2025acessibilidade} apontam que a falta de suporte institucional e as dificuldades enfrentadas pelos docentes prejudicam o acesso ao conhecimento de estudantes com deficiência intelectual no ensino fundamental. Isso mostra que garantir a presença do aluno na sala de aula não significa, necessariamente, oferecer condições adequadas para sua aprendizagem.

Diante desse cenário, a utilização de tecnologias digitais pode ajudar os professores a lidar com parte dessas dificuldades. Entre essas tecnologias, a Inteligência Artificial (IA) tem apresentado possibilidades de aplicação na educação, principalmente na criação e adaptação de conteúdos. A Inteligência Artificial Generativa permite produzir textos, exercícios e outros materiais a partir de orientações fornecidas pelo usuário. \citeonline{engaja2026} aborda a utilização de modelos de linguagem como ferramentas de apoio aos professores no ensino inclusivo, mostrando possibilidades de aplicação dessa tecnologia na elaboração de recursos educacionais.

Apesar dessas possibilidades, a utilização da Inteligência Artificial na educação exige cuidados. Os conteúdos gerados podem apresentar erros, informações inadequadas ou atividades que não correspondam às necessidades dos alunos. Por esse motivo, a tecnologia deve funcionar como uma ferramenta de apoio, mantendo o professor responsável pela avaliação e adaptação dos materiais antes de utilizá-los em sala de aula. Dessa maneira, é possível aproveitar os recursos oferecidos pela IA sem substituir o conhecimento e a experiência do profissional da educação.

Para auxiliar os professores na preparação de atividades adaptadas, foi desenvolvido o IncluiEdu, um sistema web voltado ao ensino fundamental que utiliza Inteligência Artificial Generativa na criação de materiais pedagógicos personalizados para estudantes neurodivergentes. A plataforma permite cadastrar turmas e alunos, registrar observações pedagógicas, acompanhar o desenvolvimento dos estudantes e gerar atividades de acordo com suas necessidades. Os conteúdos produzidos pela IA são armazenados inicialmente como rascunhos, permitindo que o professor faça as alterações necessárias antes de utilizá-los.

O IncluiEdu foi desenvolvido com a proposta de facilitar a organização das informações pedagógicas e apoiar a elaboração de atividades adaptadas. A ferramenta busca oferecer aos professores uma alternativa para o planejamento de materiais personalizados, especialmente em situações nas quais o tempo e os recursos disponíveis são limitados. A decisão sobre a utilização das atividades permanece sob responsabilidade do docente, que pode avaliar os conteúdos e adequá-los às necessidades de cada estudante.

\section{PROBLEMA}

A problemática central deste trabalho reside na lacuna entre o direito garantido por lei e a realidade prática do professor do ensino fundamental na rede pública. Em uma mesma turma, o docente pode trabalhar com crianças que apresentam ritmos de aprendizagem, dificuldades cognitivas e níveis de autonomia distintos (perfis neurodivergentes). Para atender a essas diferenças, é imprescindível preparar atividades e utilizar estratégias que considerem as particularidades de cada aluno.

Contudo, conforme destacam \citeonline{oliveira2025escolas}, a ausência de materiais pedagógicos adequados e de formação continuada está entre as principais dificuldades enfrentadas pelas escolas regulares na implementação da educação inclusiva. O tempo disponível para o planejamento das aulas e a falta de recursos dificultam esse processo de individualização do ensino.

Nesse contexto, embora a Inteligência Artificial Generativa se apresente como uma alternativa promissora para a criação de conteúdos, o uso dessas ferramentas exige cautela. Os conteúdos gerados podem conter erros ou não se alinhar às necessidades específicas do estudante.

Diante disso, o problema que este trabalho se propõe a responder é: Como apoiar o professor do ensino fundamental na criação de atividades pedagógicas adaptadas para alunos neurodivergentes, de forma a reduzir sua sobrecarga de trabalho, mas garantindo que o docente mantenha o controle pedagógico e a validação ética dos materiais gerados?

\section{OBJETIVOS}

\subsection{OBJETIVO GERAL}
Desenvolver o IncluiEdu, um sistema web fundamentado em Inteligência Artificial Generativa, destinado a apoiar professores do ensino fundamental na criação de atividades pedagógicas adaptadas para alunos neurodivergentes, de modo a otimizar o tempo de planejamento docente e assegurar que a validação pedagógica e ética dos materiais permaneça sob controle do professor.
\subsection{OBJETIVOS ESPECÍFICOS}

\begin{itemize}
    \item Identificar, por meio da literatura, os principais desafios enfrentados por professores do ensino fundamental público na criação de atividades adaptadas para alunos neurodivergentes.
    \item Definir os requisitos funcionais e não funcionais do sistema, com foco na segurança dos dados (LGPD) e no controle de acesso aos registros pedagógicos.
    \item Projetar a arquitetura e o modelo de banco de dados relacional para o armazenamento de turmas, alunos e histórico de observações pedagógicas.
    \item Desenvolver uma interface web inclusiva, aplicando princípios de neuroinclusão, como a paleta de cores pastéis e a sinalização cromática (badges), para a redução da carga cognitiva do professor.
    \item Implementar a integração com a Inteligência Artificial Generativa (via Cloudflare Workers AI) para a criação de atividades pedagógicas personalizadas.
    \item Implementar um recurso de geração de imagens pedagógicas por meio de Inteligência Artificial Generativa, visando complementar visualmente as atividades adaptadas produzidas pelo sistema.
    \item Estabelecer um mecanismo de segurança pedagógica no qual as atividades geradas pela IA sejam armazenadas obrigatoriamente como rascunhos, exigindo a revisão e validação final do professor.
    \item Validar o sistema por meio de testes manuais de usabilidade e funcionamento dos fluxos de cadastro, acompanhamento e geração de conteúdo.
\end{itemize}

\section{JUSTIFICATIVA}

A relevância deste trabalho fundamenta-se em três pilares principais: a urgência da inclusão educacional no Brasil, a necessidade de uso ético da Inteligência Artificial na educação e a redução da sobrecarga do trabalho docente.

No âmbito social e educacional, os dados do Censo Escolar de 2025 revelam que aproximadamente 2,4 milhões de estudantes com alguma deficiência estão matriculados na educação básica, sendo a deficiência intelectual a categoria mais numerosa \cite{inep2025censo}. Contudo, a realidade das escolas públicas demonstra uma lacuna entre o direito garantido pela Lei Brasileira de Inclusão \cite{brasil2015lbi} e a prática pedagógica. Apenas 6,4\% dos professores regentes possuem formação continuada em Educação Especial \cite{irm2026panorama}, o que, somado à falta de materiais adaptados, resulta em uma inclusão que muitas vezes se limita à socialização do aluno, sem garantir a aprendizagem efetiva \cite{oliveira2025escolas}. 

No âmbito tecnológico, a Inteligência Artificial Generativa tem se destacado como uma ferramenta promissora para a criação de conteúdo. No entanto, seu uso na educação exige cautela, pois os modelos de linguagem podem apresentar erros ou vieses. Este trabalho justifica-se por propor uma aplicação da IA pautada na ética e na segurança pedagógica \cite{engaja2026}. O diferencial do IncluiEdu é manter o professor no centro do processo decisório, por meio do sistema de ``rascunhos'', onde a IA atua como copiloto na elaboração de atividades, mas a validação humana é obrigatória. Além disso, o sistema inova ao aplicar princípios de design neuroinclusivo \cite{bertoletti2026light,kozybayev2026cognitive} na interface do docente, reduzindo a fadiga visual e a carga cognitiva.

Por fim, no âmbito prático e acadêmico, o IncluiEdu justifica-se por otimizar o tempo de planejamento do professor. Ao automatizar a criação de atividades adaptadas com base nas observações pedagógicas do próprio docente, o sistema mitiga o problema da sobrecarga de trabalho, democratizando o acesso a materiais personalizados e contribuindo diretamente para uma educação mais equitativa e verdadeiramente inclusiva.
